% Generated from project outputs. Do not edit by hand.
\begin{table}[!htbp]
\centering
\begin{threeparttable}
\caption{Labor income by subsample: 2018 moments}
\label{tab:annual-2018-y1-moments}
\small
\setlength{\tabcolsep}{4pt}
\renewcommand{\arraystretch}{1.15}
\begin{tabularx}{\linewidth}{@{}Xrrrrr@{}}
\toprule
Sample & N & Mean & SD & Min & Max \\
\midrule
General & 94 & 2,690.5 & 2,183.9 & 279.4 & 16,913.6 \\
Men & 94 & 3,014.2 & 2,762.1 & 199.3 & 22,755.4 \\
Women & 94 & 2,106.7 & 1,425.2 & 241.0 & 7,275.7 \\
Urban & 94 & 3,053.1 & 2,262.4 & 283.2 & 18,300.5 \\
Rural & 93 & 1,796.0 & 1,075.6 & 241.7 & 5,390.6 \\
Indigenous & 93 & 1,872.0 & 1,178.3 & 211.9 & 5,507.4 \\
Non-Indigenous & 94 & 2,899.7 & 2,330.6 & 193.0 & 18,863.6 \\
Bottom 99 percent & 94 & 2,397.5 & 1,406.6 & 279.4 & 5,681.6 \\
Bottom 90 percent & 94 & 1,909.4 & 918.5 & 299.0 & 4,079.3 \\
Bottom 50 percent & 94 & 942.7 & 426.4 & 120.3 & 2,668.3 \\
\bottomrule
\end{tabularx}
\begin{tablenotes}[flushleft]
\footnotesize
\item Monthly quetzales. Unweighted distribution of survey-weighted cohort means before transition matching. N counts cohort cells, not individuals. SD is dispersion, not a standard error. Subsamples overlap. Cumulative income definitions and treatment of missing components follow the merging scripts.
\end{tablenotes}
\end{threeparttable}
\end{table}
